\section{Trabalhos Relacionados}
\label{sec:trabalhos}

Esta seção discute cinco estudos empíricos da engenharia de \emph{software} que tratam da qualidade de código assistido por IA, da corretude de artefatos gerados por modelos de linguagem ou da interpretação de métricas de acoplamento em Java. Para cada trabalho, apresentam-se o objetivo e o método, os resultados pertinentes a este TCC e a relação com a proposta --- complementar, comparar ou reutilizar o mesmo tipo de evidência. Os três primeiros artigos foram publicados em 2023 e estão diretamente ligados ao recorte da Era da IA assistida.

A qualidade do código produzido por assistentes foi avaliada de forma experimental com o conjunto HumanEval, comparando GitHub Copilot, Amazon CodeWhisperer e ChatGPT quanto a validade, correção, segurança, confiabilidade e manutenibilidade \cite{yetistiren:23}. O método consistiu em gerar soluções a partir de enunciados e inspecionar o código resultante com métricas de qualidade e critérios de aceitação dos testes do \emph{benchmark}. Os resultados indicam taxas de correção de 65,2\% para o ChatGPT, 46,3\% para o Copilot e 31,1\% para o CodeWhisperer, além de dívida técnica média associada a \emph{code smells} da ordem de minutos por trecho. Esse estudo é complementar ao presente trabalho: enquanto avalia trechos gerados em tarefas isoladas, esta pesquisa observa a estrutura de repositórios Spring Boot reais ao longo do tempo, permitindo contrastar o período pré-IA com a Era da IA assistida.

A capacidade do GitHub Copilot como programador-par foi examinada em dois conjuntos de tarefas: problemas algorítmicos fundamentais e um corpus de soluções humanas, corretas e defeituosas \cite{dakhel:23}. O método combinou geração repetida de soluções, verificação de funcionalidade e comparação com implementações de programadores. Os resultados mostram que o Copilot produz soluções para a maioria dos problemas, porém parte delas é defeituosa ou não reproduzível; a taxa de correção humana superou a do assistente, embora o código gerado com falhas tenda a exigir menor esforço de reparo. Os autores concluem que a ferramenta pode ser um ativo para desenvolvedores experientes e um passivo para novatos. O presente trabalho não replica o desenho de tarefas algorítmicas; em vez disso, complementa essa evidência ao perguntar se a adoção massiva de assistentes, a partir de 2022, coincide com mudanças nas métricas estruturais de projetos Java/Spring Boot publicados no GitHub.

A correção funcional de código gerado por modelos de linguagem, incluindo o ChatGPT, foi avaliada com o \emph{framework} EvalPlus, que amplia o HumanEval com uma quantidade substancialmente maior de casos de teste obtidos por geração automática e mutação \cite{liu:23}. O método consistiu em reexecutar 26 modelos populares sobre HumanEval e HumanEval+, medindo \emph{pass@k}. Os resultados mostram quedas de até 19,3\%--28,9\% em \emph{pass@k} quando os testes adicionais são considerados, e alteração no ranqueamento de modelos --- inclusive casos em que o ChatGPT deixa de superar alternativas de código aberto. A relação com este trabalho é de complementaridade: os resultados de \cite{liu:23} demonstram que a qualidade funcional do código gerado por IA é frequentemente superestimada em \emph{benchmarks} curtos; o presente estudo desloca o foco da correção de funções sintéticas para a qualidade estrutural de sistemas Spring Boot completos, em coortes históricas.

O impacto da injeção de dependências sobre a manutenibilidade em projetos Java foi investigado com duas métricas ponderadas --- DCE e DCBO --- implementadas no CKJM-Analyzer, extensão da ferramenta CKJM \cite{sun:22}. O método aplicou análise estática a um conjunto de projetos \emph{open source}, distinguindo acoplamento ``suave'' (via injeção) de acoplamento ``rígido'' (instanciação direta). Os resultados indicam que a injeção de dependências reduz o acoplamento efetivo quando este é ponderado, o que métricas clássicas de CBO tendem a não capturar integralmente. O presente trabalho compartilha o objeto de estudo --- código Java com forte presença de injeção, típica do Spring Boot ---, porém não propõe métricas novas: utiliza a suíte CK de forma consistente entre épocas para observar evolução estrutural, reconhecendo a limitação apontada por \cite{sun:22} na interpretação do acoplamento.

O efeito do GitHub Copilot sobre a produtividade de desenvolvedores profissionais foi medido em experimento controlado, com tarefa de implementação cronometrada e grupo de controle sem o assistente \cite{peng:23}. O método alocou participantes de forma aleatória e comparou tempo de conclusão e taxa de término. Os resultados apontam redução substancial do tempo necessário para completar a tarefa (cerca de 55\% mais rápido no grupo com Copilot) e maior taxa de conclusão. Essa evidência fundamenta a delimitação da quarta época histórica (a partir de 2022) como Era da IA assistida. A diferença em relação a este trabalho é o construto medido: o experimento de \cite{peng:23} avalia desempenho individual em uma tarefa; aqui se avaliam métricas estruturais de repositórios públicos, sem reivindicar causalidade direta entre o uso de um assistente específico e a qualidade de cada projeto.

Em conjunto, os estudos revisados mostram que assistentes de IA alteram produtividade e qualidade de trechos gerados, e que métricas clássicas de acoplamento exigem cautela em sistemas com injeção de dependências. Nenhum deles, contudo, descreve a trajetória longitudinal da qualidade estrutural de projetos Spring Boot no GitHub. Essa lacuna é o ponto de partida da metodologia apresentada a seguir.
